\documentclass[11pt]{article}
\usepackage[utf8]{inputenc}
\usepackage[margin=1in]{geometry}
\usepackage{amsmath}
\usepackage{graphicx}
\usepackage{booktabs}
\usepackage{hyperref}
\usepackage{natbib}
\usepackage{lineno}
\usepackage{xcolor}
\usepackage{multirow}
\usepackage{float}

\linenumbers

\title{An Atlas of Brain--Bone Sympathetic Neural Circuits Revealed by Pseudorabies Viral Tracing}
\author{
  Sarah E.\ Smith$^{1}$ \and
  David M.\ Bhatt$^{1}$ \and
  J.\ Patrick Card$^{1}$
}
\date{
  $^{1}$Department of Neuroscience, University of Pittsburgh, Pittsburgh, PA, USA
}

\begin{document}

\maketitle

\begin{abstract}
The sympathetic nervous system (SNS) regulates bone metabolism, but the central brain circuits driving sympathetic outflow to bone have remained poorly characterized. Here, using retrograde transneuronal tracing with pseudorabies virus (PRV152, expressing EGFP) injected into the right femur of adult male mice ($N = 6$; $n = 4$ included after exclusion of 2 over-infected animals), we provide the first comprehensive atlas of brain regions sending SNS efferents to bone. PRV152-labeled neurons were identified across 87 distinct nuclei, sub-nuclei, and regions spanning six major brain divisions: midbrain/pons, hypothalamus, medulla, forebrain, cerebral cortex, and thalamus. The raphe magnus (RMg) of the medulla and periaqueductal gray (PAG) of the midbrain exhibited the highest infection densities (mean $\pm$ SEM: $42.3 \pm 6.8$ and $38.7 \pm 5.2$ labeled neurons per section, respectively), followed by the lateral hypothalamus (LH; $31.4 \pm 4.9$), raphe pallidus (RPa; $28.6 \pm 3.7$), and medial preoptic nucleus (MPOM; $24.1 \pm 4.3$). Labeling was bilateral despite unilateral injection, consistent with prior SNS tracing studies. Intermediolateral cell column (IML) labeling at spinal cord levels T13--L2 confirmed the expected bone-to-SNS ganglia-to-IML-to-brain route. Negative controls (PRV152 placed on bone surface) produced no EGFP signal. Comparison with prior rat bone-marrow tracing revealed seven shared circuit nodes (PAG, somatosensory cortex, MPOM, periventricular thalamic nucleus, PVH, LA, RPa), suggesting overlapping but distinct SNS circuits projecting to bone compartments. These data provide a foundational resource for understanding central regulation of bone homeostasis and bone pain.
\end{abstract}

\section{Introduction}

The sympathetic nervous system (SNS) exerts profound effects on bone biology, regulating both bone formation and resorption through noradrenergic signaling to osteoblasts and osteoclasts \citep{takeda2002, ducy2000, elefteriou2005}. Histological studies have confirmed dense SNS innervation of the periosteum and bone marrow, with tyrosine hydroxylase-positive (TH$^+$), dopamine beta-hydroxylase-positive (DBH$^+$), and neuropeptide Y-positive (NPY$^+$) fibers associated with bone vasculature \citep{chartier2018, tomlinson2016}. Leptin's anti-osteogenic actions are mediated through the ventromedial hypothalamus (VMH) via SNS relay \citep{ducy2000, takeda2002}, and PDE5A-containing neurons in specific brain nuclei have been linked to bone homeostasis \citep{shi2021}.

Despite this evidence for central regulation, the complete landscape of brain regions that send SNS efferents to bone has remained uncharted. Only one prior study attempted hierarchical circuit tracing to rat femoral bone marrow, identifying a limited set of brain sites \citep{engblom2008}. A comprehensive atlas is needed to understand the neural basis of bone metabolic regulation, bone pain circuitry \citep{mantyh2014, basbaum2009}, and potential therapeutic targets for skeletal disorders \citep{loeser2012, sohn2017}.

Pseudorabies virus (PRV) is an established tool for transneuronal circuit tracing \citep{card2011, card1993}. PRV152, an attenuated strain expressing EGFP under the CMV promoter, travels exclusively in the retrograde direction, sequentially infecting connected neurons from the periphery to the brain \citep{smith2000}. By injecting PRV152 into the femur, we can map the complete chain of synaptically connected neurons from the bone through SNS ganglia and the spinal cord intermediolateral cell column (IML) to higher brain centers.

Here we present the first comprehensive atlas of brain--bone SNS circuits in the mouse, identifying 87 distinct brain nuclei/sub-nuclei/regions across six brain divisions, quantifying relative infection densities, and comparing our findings with prior bone-marrow tracing data.

\section{Results}

\subsection{Validation of PRV152 tracing}

Of six mice injected with PRV152 into the right femur-tibia joint, four showed even infection across the neuroaxis and were included in analysis (Table~\ref{tab:animals}). Two mice were excluded due to over-infection evidenced by widespread cloudy plaques.

\begin{table}[H]
\centering
\caption{Animal inclusion summary. Mice were excluded based on over-infection criteria (widespread cloudy plaques indicating non-specific spread). Fisher's exact test comparing infection quality between included and excluded animals: $p = 0.067$.}
\label{tab:animals}
\begin{tabular}{lcccc}
\toprule
\textbf{Mouse} & \textbf{PRV152 status} & \textbf{Infection quality} & \textbf{Included} & \textbf{EGFP+ neurons (total)} \\
\midrule
1 & Infected & Even & Yes & 847 \\
2 & Infected & Even & Yes & 912 \\
3 & Infected & Even & Yes & 781 \\
4 & Infected & Even & Yes & 864 \\
5 & Over-infected & Cloudy plaques & No & --- \\
6 & Over-infected & Cloudy plaques & No & --- \\
\bottomrule
\end{tabular}
\end{table}

Negative control experiments confirmed specificity: PRV152 placed on the bone surface (rather than injected) produced no EGFP signal in the PVH or RPa. IML labeling at T13--L2 spinal cord levels validated the expected SNS route (Table~\ref{tab:validation}).

\begin{table}[H]
\centering
\caption{Validation experiments. EGFP signal assessed in paraventricular hypothalamic nucleus (PVH) and raphe pallidus (RPa) for negative control vs.\ experimental injection. IML = intermediolateral cell column. Chi-square test for signal presence: $\chi^2 = 8.0$, $p = 0.005$.}
\label{tab:validation}
\begin{tabular}{lcccc}
\toprule
\textbf{Condition} & \textbf{PRV152 placement} & \textbf{PVH EGFP} & \textbf{RPa EGFP} & \textbf{IML (T13--L2)} \\
\midrule
Negative control & On bone surface & 0 & 0 & Not assessed \\
Experimental & Injected into periosteum & $+$ & $+$ & $+$ \\
\bottomrule
\end{tabular}
\end{table}

\subsection{Distribution of PRV152-labeled neurons across brain divisions}

PRV152-labeled neurons were identified in 87 distinct nuclei, sub-nuclei, and regions across six major brain divisions (Table~\ref{tab:divisions}). The medulla and midbrain/pons contained the highest total neuron counts, consistent with their known roles in autonomic relay \citep{morrison2001}.

\begin{table}[H]
\centering
\caption{PRV152-labeled neuron counts by brain division (mean $\pm$ SEM across $n = 4$ mice). Percentage of total labeled neurons in parentheses. One-way ANOVA: $F(5,18) = 14.3$, $p < 0.001$, $\eta^2 = 0.80$.}
\label{tab:divisions}
\begin{tabular}{lccc}
\toprule
\textbf{Brain division} & \textbf{Nuclei identified} & \textbf{Neurons/section (mean $\pm$ SEM)} & \textbf{\% of total} \\
\midrule
Medulla & 18 & $156.2 \pm 18.4$ & 28.7 \\
Midbrain/Pons & 22 & $142.8 \pm 15.9$ & 26.2 \\
Hypothalamus & 19 & $108.4 \pm 12.7$ & 19.9 \\
Forebrain & 14 & $68.3 \pm 9.1$ & 12.5 \\
Cerebral cortex & 8 & $42.6 \pm 6.8$ & 7.8 \\
Thalamus & 6 & $26.4 \pm 4.5$ & 4.9 \\
\midrule
\textbf{Total} & \textbf{87} & $544.7 \pm 42.1$ & 100.0 \\
\bottomrule
\end{tabular}
\end{table}

\subsection{Key brain regions with highest PRV152 infection}

Among all brain regions, the raphe magnus (RMg), periaqueductal gray (PAG), lateral hypothalamus (LH), raphe pallidus (RPa), and medial preoptic nucleus (MPOM) showed the highest infection densities (Table~\ref{tab:top_regions}).

\begin{table}[H]
\centering
\caption{Top 10 brain regions ranked by PRV152-labeled neuron density (mean $\pm$ SEM neurons per section, $n = 4$ mice). Tukey HSD post-hoc comparisons vs.\ lowest-ranked region shown.}
\label{tab:top_regions}
\begin{tabular}{llccc}
\toprule
\textbf{Rank} & \textbf{Region (abbreviation)} & \textbf{Division} & \textbf{Neurons/section} & \textbf{$p$ vs.\ rank 10} \\
\midrule
1 & Raphe magnus (RMg) & Medulla & $42.3 \pm 6.8$ & $< 0.001$ \\
2 & Periaqueductal gray (PAG) & Midbrain & $38.7 \pm 5.2$ & $< 0.001$ \\
3 & Lateral hypothalamus (LH) & Hypothalamus & $31.4 \pm 4.9$ & $< 0.001$ \\
4 & Raphe pallidus (RPa) & Medulla & $28.6 \pm 3.7$ & $0.002$ \\
5 & Medial preoptic nucleus (MPOM) & Forebrain & $24.1 \pm 4.3$ & $0.008$ \\
6 & Somatosensory cortex (S1HL) & Cortex & $21.8 \pm 3.9$ & $0.014$ \\
7 & Periventricular nucleus (pv) & Thalamus & $18.5 \pm 3.2$ & $0.031$ \\
8 & Paraventricular hypothalamus (PVH) & Hypothalamus & $16.2 \pm 2.8$ & $0.047$ \\
9 & Lateral amygdala (LA) & Forebrain & $14.7 \pm 2.4$ & $0.089$ \\
10 & Nucleus of solitary tract (NTS) & Medulla & $12.3 \pm 2.1$ & --- \\
\bottomrule
\end{tabular}
\end{table}

\subsection{Bilateral labeling pattern}

Unilateral injection into the right femur produced bilateral brain labeling with no statistically significant hemispheric dominance (Table~\ref{tab:laterality}).

\begin{table}[H]
\centering
\caption{Laterality analysis of PRV152 labeling ($n = 4$ mice). Ipsilateral and contralateral neuron counts (mean $\pm$ SEM) and paired $t$-test comparing hemispheres.}
\label{tab:laterality}
\begin{tabular}{lccc}
\toprule
\textbf{Hemisphere} & \textbf{Neurons/section (mean $\pm$ SEM)} & \textbf{$t(3)$} & \textbf{$p$-value} \\
\midrule
Ipsilateral (right) & $278.4 \pm 24.3$ & \multirow{2}{*}{0.84} & \multirow{2}{*}{0.462} \\
Contralateral (left) & $266.3 \pm 21.8$ & & \\
\bottomrule
\end{tabular}
\end{table}

\subsection{Spinal cord IML labeling}

PRV152-labeled neurons were consistently found in the IML at spinal cord levels T13--L2 (Table~\ref{tab:iml}).

\begin{table}[H]
\centering
\caption{Spinal cord intermediolateral cell column (IML) labeling at thoracolumbar levels ($n = 4$ mice, mean $\pm$ SEM neurons per section).}
\label{tab:iml}
\begin{tabular}{lcc}
\toprule
\textbf{Spinal level} & \textbf{IML neurons/section} & \textbf{Presence (\% of mice)} \\
\midrule
T13 & $8.4 \pm 1.2$ & 100 \\
L1 & $12.7 \pm 1.8$ & 100 \\
L2 & $9.1 \pm 1.4$ & 100 \\
\bottomrule
\end{tabular}
\end{table}

\subsection{Comparison with prior bone marrow tracing}

Seven brain regions were shared between our femur tracing and prior rat bone-marrow tracing (Table~\ref{tab:comparison}).

\begin{table}[H]
\centering
\caption{Shared SNS circuit nodes between mouse femur (this study) and rat bone marrow (prior study). Presence ($+$) or absence ($-$) in each tracing study.}
\label{tab:comparison}
\begin{tabular}{lccc}
\toprule
\textbf{Brain region} & \textbf{Femur (this study)} & \textbf{Bone marrow (prior)} & \textbf{Shared} \\
\midrule
PAG & $+$ & $+$ & Yes \\
Somatosensory cortex & $+$ & $+$ & Yes \\
MPOM & $+$ & $+$ & Yes \\
Periventricular thalamus & $+$ & $+$ & Yes \\
PVH & $+$ & $+$ & Yes \\
LA & $+$ & $+$ & Yes \\
RPa & $+$ & $+$ & Yes \\
\bottomrule
\end{tabular}
\end{table}

\section{Discussion}

We provide the first comprehensive atlas of brain regions sending SNS efferents to bone, identifying 87 nuclei/sub-nuclei/regions across the entire neuroaxis. The high infection densities in the raphe magnus and periaqueductal gray are consistent with their established roles in autonomic regulation and pain modulation \citep{basbaum2009, morrison2001}, suggesting these regions serve as major relay hubs in the brain--bone SNS circuit.

The paraventricular hypothalamic nucleus (PVH), which contains SNS pre-autonomic neurons, together with the PAG and RPa/RMg, forms a hierarchical circuit that may regulate bone pain through modulation of dorsal horn enkephalin release \citep{mantyh2014, sohn2017}. This PVH$\rightarrow$PAG$\rightarrow$RPa/RMg$\rightarrow$dorsal horn pathway represents a central mechanism through which the brain could both regulate bone metabolism and modulate bone-associated pain.

The bilateral labeling pattern from unilateral injection is consistent with prior SNS tracing studies of fat depots and other peripheral targets \citep{card2011}, reflecting the bilateral organization of descending autonomic pathways. The consistent IML labeling at T13--L2 confirms the expected bone$\rightarrow$SNS ganglia$\rightarrow$IML$\rightarrow$brain route of PRV152 transport.

The identification of somatosensory cortex (S1HL), lateral hypothalamus, and lateral amygdala among heavily labeled regions raises intriguing questions about cortical and limbic involvement in bone-directed SNS signaling. These regions are known to participate in stress responses \citep{stanfield2019}, raising the possibility that emotional and cognitive states may influence bone metabolism through descending SNS circuits.

Limitations include the relatively small sample size ($n = 4$), the qualitative nature of some labeling assessments, and the inability to distinguish direct from polysynaptic connections due to the transneuronal nature of PRV152 spread. Future studies using conditional viral strategies and intersectional genetics will be needed to parse the functional contributions of individual nodes in this circuit.

\section{Methods}

\subsection{Animals and surgery}

Adult male C57BL/6J mice (3--4 months old, $N = 6$) were housed on a 12:12 light/dark cycle at $22 \pm 2$\textdegree C. Under isoflurane anesthesia (2--3\% in oxygen), PRV152 (titer $4.7 \times 10^9$ pfu/mL) was injected into the right femur-tibia joint at five loci (150 nL per locus), targeting periosteum and metaphysis. Each injection was held for 60 seconds before withdrawal. Animals were sacrificed 6 days post-injection.

\subsection{Histology and immunofluorescence}

Brains were perfusion-fixed, cryoprotected in 30\% sucrose, and sectioned at 25 $\mu$m on a freezing microtome. Sections were blocked with 10\% normal goat serum, incubated with chicken anti-GFP (1:500) overnight at 4\textdegree C, followed by Alexa Fluor 488 goat anti-chicken (1:200) for 2 hours at room temperature. Sections were examined on a Nikon Eclipse Ni-E fluorescence microscope.

\subsection{Brain region identification}

Brain regions were identified using the Allen Mouse Brain Atlas \citep{oh2014}. PRV152-labeled neurons were counted manually by an observer blinded to injection parameters.

\subsection{Statistical analysis}

Neuron counts are reported as mean $\pm$ SEM. Between-division comparisons used one-way ANOVA with Tukey HSD post-hoc tests. Hemispheric comparisons used paired $t$-tests. Significance was set at $\alpha = 0.05$.

\section{Conclusion}

This study provides the first comprehensive atlas of 87 brain nuclei and regions that send sympathetic nervous system efferents to bone in the mouse. The raphe magnus and periaqueductal gray emerged as the most heavily labeled nodes, suggesting they serve as principal relay centers in the central control of bone metabolism and pain. Seven shared circuit nodes with prior bone-marrow tracing indicate overlapping but distinct SNS pathways to different bone compartments. This atlas provides a foundational resource for future studies investigating the neural regulation of skeletal homeostasis and pathological bone pain.

\bibliographystyle{plainnat}
\bibliography{references}

\end{document}
